\begin{corollary}[7.4.5]
\label{IV.2.s7.4.5-fr}
Supongamos que la propiedad $\boldsymbol P$ satisface las condiciones
$(\mathrm P_{\mathrm I})$, $(\mathrm P_{\mathrm{II}})$ y
$(\mathrm P_{\mathrm{III}})$. Sea $A$ un anillo noetheriano. Las
condiciones siguientes son equivalentes:
\begin{enumerate}
\item[\emph{a)}] Para todo ideal primo $\mathfrak p$ de $A$,
$A_{\mathfrak p}$ es un $\boldsymbol P$-anillo.
\item[\emph{b)}] Para todo ideal maximal $\mathfrak m$ de $A$,
$A_{\mathfrak m}$ es un $\boldsymbol P$-anillo.
\end{enumerate}
Si además $\boldsymbol P$ satisface $(\mathrm P_{\mathrm{IV}})$, las
propiedades \textup{a)} y \textup{b)} son también equivalentes a:
\begin{enumerate}
\item[\emph{c)}] Para toda $A$-álgebra de tipo finito $B$ y todo ideal
primo $\mathfrak q$ de $B$, $B_{\mathfrak q}$ es un
$\boldsymbol P$-anillo.
\end{enumerate}
\end{corollary}

\begin{proof}
La equivalencia de \emph{a)} y \emph{b)} resulta de (7.4.4, (i)), y la de \emph{a)} y \emph{c)} de
(7.4.4, (ii)).

Cuando se cumple la condición \emph{a)} de (7.4.5), se dice que $A$ es un
$\boldsymbol P$-\emph{anillo}; para los anillos semilocales noetherianos, esta
definición coincide con la de (7.3.13) cuando se satisfacen las
condiciones $(\mathrm P_{\mathrm I})$, $(\mathrm P_{\mathrm{II}})$ y
$(\mathrm P_{\mathrm{III}})$. El anillo $\mathbf Z$ es un
$\boldsymbol P$-anillo (7.3.19, (iv)); todo anillo local completo es un
$\boldsymbol P$-anillo.
\end{proof}
